% Iter 23 — Pillar 3 elevation: phase diagram, T_critical, Pareto frontier,
% anti-herding δ_div injection simulation.
% (M. M3 autonomous, July 2026.)

\subsubsection{Pillar 3 elevation: phase diagram, $T_\text{crit}$, and the anti-herding lever}
\label{sec:pillar3-iter23}

\paragraph{Setup.}
We lift the Pillar-3 question~\cite{tan2025scalingrl} ``does $G\!\approx\!32$ confer a meaningful
benefit over $G\!=\!4$?'' from point estimates to a four-axis decomposition:
\begin{enumerate}\itemsep1pt
\item a continuous $T_\text{crit}$ estimate (when does $G\!=\!4$ cease to be
      equivalent to $G\!=\!32$ on a paired metric);
\item an illustrative accuracy$(T, G)$ phase diagram reconstructed for
      Qwen3-8B / GSM8K;
\item the within-reconstruction argmax $G^\star(T)$ frontier;
\item an analytical simulation of how much diversity amplification
      $\delta_\text{div}$ is required at $G\!=\!4$ to recover the
      $G\!=\!32$ contrastive yield.
\end{enumerate}

\paragraph{$T_\text{crit}$ for $G=4$ vs $G=32$.}
Fitting $|\Delta(T)| = a \cdot (1 - e^{-T/\tau})$ to four reconstructed
budget points
yields $\hat a=0.30$, $\hat\tau \!\approx\! 12.1$M tokens.
The 5pp-equivalence breaking threshold falls at
$T_\text{crit}^{5\%}\!=\!2.21$M tokens
(bootstrap 95\% CI: $[1.29\text{M}, 4.77\text{M}]$).  At $T\!=\!1$M, the
canonical comparison holds ($|\Delta|\!\le\!1$pp); at $T\!=\!4$M it breaks.
Numerically, this reconstruction places the boundary between
``Wu-equivalence operational'' and ``$G$ matters'' between $1$M and $2$M
tokens.  It is a hypothesis for a matched-budget Qwen3-8B/GSM8K sweep, not a
measured learning-curve transition.

\paragraph{Phase diagram and Pareto frontier.}
Table~\ref{tab:p3-pareto} shows the per-budget argmax inside the reconstructed
grid.  It must not be read as a deployment optimum.

\begin{table}[h!]
\centering
\caption{Within-reconstruction $G^\star(T)$ on Qwen3-8B / GSM8K, with
conditional bootstrap 95\% intervals.}
\label{tab:p3-pareto}
\small
\begin{tabular}{rccc}
\hline
$T$ (tokens) & $G^\star(T)$ & acc($G^\star$) & CI \\ \hline
  1,000,000 &   8 & 0.48 & [0.45, 0.51] \\
  4,000,000 &  16 & 0.69 & [0.66, 0.72] \\
 16,000,000 &  32 & 0.84 & [0.81, 0.87] \\
 64,000,000 &  32 & 0.88 & [0.85, 0.91] \\ \hline
\end{tabular}
\end{table}

The reconstructed surface suggests three regimes: an early argmax near
$G{=}8$, a shift toward $G{=}16$ at $T{=}4$M, and an argmax near $G{=}32$ at
the two largest budgets.  These are candidate operating points for the
matched-token experiment.  They neither establish statistical
indistinguishability of $G{=}32$ and $G{=}64$ nor resolve the Wu et al.\ result
outside the setting that paper actually tested.

\paragraph{Crossover matrix.}
Table~\ref{tab:p3-crossings} applies paired-bootstrap calculations to the
reconstructed grid.  Conditional on that grid, the largest contrast is the
$G{=}4$ versus $G{=}32$ comparison at $T{=}64$M.  Because the inputs are not
direct matched-budget runs, these intervals quantify reconstruction
uncertainty only.

\begin{table}[h!]
\centering
\caption{Conditional $\Delta_\text{pp}=\text{acc}(G_b)-\text{acc}(G_a)$ in the
reconstructed Qwen3-8B/GSM8K grid; $T_\text{cross}$ is an interpolated
threshold, not an observed training transition.}
\label{tab:p3-crossings}
\small
\IfFileExists{../experiments/results/group_size_iter23_crossings.tsv_typesetter}{\input{../experiments/results/group_size_iter23_crossings.tsv_typesetter}}{}
\end{table}

\paragraph{Anti-herding $\delta_\text{div}$ injection (frontier synthesis).}
The frontier synthesis (Round 2) crystallises the question:
\emph{``can we inflate $\delta_\text{div}$ during decoding to achieve $G\!=\!32$ yield at
$G\!=\!4$ FLOP cost?''}
Model ZVF$(G, p_\text{eff}, \delta_\text{div}) = \max(0, p^G + (1-p)^G - \delta_\text{div})$.
Solving $\text{ZVF}(4, p, \delta') = \text{ZVF}(32, p, \delta_\text{base}=0.18)$ across
$p_\text{eff}\!\in\![0.60, 0.98]$ we obtain:\newline
\textbf{feasible} ($\delta'\!\le\!0.23$): $p_\text{eff}\!\le\!0.68$ ($\sim\!23\%$ of frontier);\newline
\textbf{stretch via decoding penalty} ($0.23\!<\!\delta'\!\le\!0.45$): $p_\text{eff}\!\in\!(0.68, 0.83]$ ($\sim\!33\%$);\newline
\textbf{infeasible} ($\delta'\!>\!0.45$): $p_\text{eff}\!>\!0.83$ ($\sim\!44\%$).
Under this proxy, the Pillar-2 diversity value ($\sim\!0.23$) does not close
the modeled $G{=}4\mid G{=}32$ contrastive-yield gap in the high-$p$ region.
This is a mechanism hypothesis for why a larger group might help late in
training; it is not evidence that $G{=}32$ is empirically Pareto-optimal.

\paragraph{Figure.}
Figure~\ref{fig:p3-iter23} packs the four artifacts above into a single 4-panel view:\ (a)
$T_\text{crit}$ sweep with the 5pp-equivalence line;\ (b) per-pair Δ-pp at
$T\!=\!64$M;\ (c) the accuracy$(T, G)$ phase diagram with iso-acc contours;\ (d) the
$\delta_\text{div}$-needs curve across the learning frontier.

\begin{figure}[ht]
\centering
\includegraphics[width=\textwidth]{figures/group_size_iter23.pdf}
\caption{Pillar-3 elevation panel. \textbf{(a)} $\!|\Delta(G\!=\!4\!-\!G\!=\!32)|$ vs.
training tokens (paired bootstrap); the dashed vertical lines mark
$T_\text{crit}^{1\%}\!=\!0.41$M, $T_\text{crit}^{3\%}\!=\!1.27$M,
$T_\text{crit}^{5\%}\!=\!2.21$M, $T_\text{crit}^{10\%}\!=\!4.91$M.
\textbf{(b)}\ Per-pair $\Delta_\text{pp}$ at $T\!=\!64$M tokens; positive bars mark
pairs where the larger $G_b$ wins (\textit{frontier synthesis}).
\textbf{(c)}\ illustrative acc$(T,G)$ reconstruction for Qwen3-8B/GSM8K;
iso-acc contours interpolate the grid. \textbf{(d)}\ $\delta_\text{div}$ amplification required
for $\text{ZVF}(G\!=\!4)\!\to\!\text{ZVF}(G\!=\!32)$ vs.\ $p_\text{eff}$, with the
$\delta_\text{div}\!\le\!0.23$ natural ceiling and the $\delta_\text{div}\!\le\!0.45$
stretch (decoding penalty) limits.}
\label{fig:p3-iter23}
\end{figure}

\paragraph{Headline.}
The reconstruction predicts that $G{=}4$ and $G{=}32$ separate as token
budget grows, while the proxy model predicts that a bounded diversity lever
will not reproduce the larger group's contrastive yield over much of the
high-$p$ region.  The publishable claim is therefore a preregistered
interaction---group size by token budget---to be tested with direct runs,
not a universal $G{=}32$ recommendation.

\paragraph{Inputs and reproducibility.}
All four artifacts in Table~\ref{tab:p3-pareto}, Table~\ref{tab:p3-crossings},
the $\delta_\text{div}$ summary, and Figure~\ref{fig:p3-iter23} are generated by
\texttt{platform\_modal/scripts/group\_size\_iter23.py} on a single CPU.
Pre-existing inputs are
\texttt{group\_size\_g4\_vs\_g32\_broader\_scale.tsv} (Pillar-3 retention sweep)
and \texttt{group\_size\_token\_normalized.tsv} (the illustrative 5$\times$4 grid).
Total runtime $\!\approx\!5$s; total wall-clock for prior runs is unchanged.
